% =============================================================
% robustness_template.tex
% Standard ladder for empirical-applied papers.
% =============================================================

\section{Robustness}\label{sec:robustness}

We probe the robustness of our main result along five dimensions:
inference, alternative samples, alternative controls, alternative
estimators, and placebo tests. \Cref{tab:robust} reports the first
four dimensions; the placebo tests are in \Cref{tab:placebo}.

\input{tabs/table_robust.tex}

\paragraph{Inference.} Column~(2) clusters two-way (unit and year)
instead of one-way; standard errors are [SIMILAR / LARGER]. Column~(3)
reports wild cluster bootstrap p-values from
\citet{CameronGelbachMiller2008} via \texttt{boottest}, which agree
with conventional inference.

\paragraph{Alternative sample.} Column~(4) restricts to
[SUBSAMPLE: e.g.\ rural counties / pre-2015 period / non-pilot states]
to address [CONCERN]. The point estimate is [SIMILAR].

\paragraph{Alternative controls.} Column~(5) augments the baseline
with [CONTROL SET]. Adding controls changes the point estimate by
less than [X]\%, suggesting that selection on observables is not
driving the result.

\paragraph{Alternative estimators.} \Cref{tab:did_estimators} reports
the main effect estimated using TWFE, Callaway-Sant'Anna (2021),
Sun-Abraham (2021), and Borusyak-Jaravel-Spiess (2024). The four
estimators agree to within [X]\%, suggesting that effect heterogeneity
across cohorts is not driving our conclusion.

\input{tabs/table_did_estimators.tex}

\paragraph{Placebo.} \Cref{tab:placebo} reports the same specification
applied to [PLACEBO OUTCOME], which the program should not affect.
The estimated effect is [POINT ESTIMATE] (s.e.\ [SE]) -- a precise
null.

\input{tabs/table_placebo.tex}

\paragraph{Pre-trends.} \Cref{fig:event} confirms that pre-period
coefficients in the event-study specification are jointly
indistinguishable from zero ($p = [P]$).
